% ==============================================================================
% PREAMBLE.TEX - Configuración y Paquetes para Teoría de la Probabilidad
% ==============================================================================

\usepackage[spanish,es-tabla]{babel}
\usepackage[utf8]{inputenc}
\usepackage{amsmath, amsfonts, amssymb, amsthm, mathtools}
\usepackage{enumitem}
\usepackage{hyperref}
\usepackage{xcolor}
\usepackage[most]{tcolorbox}
\tcbuselibrary{theorems, skins, breakable}
\usepackage{bookmark}

% Asegura el registro de capítulos y secciones en el índice
\setcounter{tocdepth}{2}
\setcounter{secnumdepth}{2}

% ------------------------------------------------------------------------------
% Configuración de Geometría y Distribución Vertical
% ------------------------------------------------------------------------------
\usepackage[margin=2.5cm, headheight=14pt]{geometry}
\raggedbottom % Permite altura natural al final de página con cajas e imágenes

% ------------------------------------------------------------------------------
% Configuración de Encabezados y Pies de Página a Doble Cara (fancyhdr)
% ------------------------------------------------------------------------------
\usepackage{fancyhdr}
\pagestyle{fancy}
\fancyhf{} % Limpia todos los campos

% Configuración para páginas Pares (Página izquierda / Par)
\fancyhead[LE]{\small\bfseries\leftmark}  % Capítulo justificado a la izquierda
\fancyhead[RE]{\small\bfseries\thepage}   % Número de página a la derecha

% Configuración para páginas Impares (Página derecha / Impar)
\fancyhead[LO]{\small\bfseries\thepage}   % Número de página a la izquierda
\fancyhead[RO]{\small\bfseries\rightmark}  % Sección justificada a la derecha

% Formato de las etiquetas de capítulo y sección
\renewcommand{\chaptermark}[1]{\markboth{\thechapter.\ #1}{}}
\renewcommand{\sectionmark}[1]{\markright{\thesection.\ #1}}
\renewcommand{\headrulewidth}{0.4pt} % Línea del encabezado

% Estilo plano para páginas de inicio de capítulo
\fancypagestyle{plain}{
  \fancyhf{}
  \fancyfoot[C]{\thepage}
  \renewcommand{\headrulewidth}{0pt}
}

% ------------------------------------------------------------------------------
% Configuración de TikZ y solución de conflictos con Babel Spanish
% ------------------------------------------------------------------------------
\usepackage{tikz}
\usetikzlibrary{shapes.geometric, arrows.meta, positioning, fit, backgrounds, babel}

% ------------------------------------------------------------------------------
% Configuración de Hipervínculos
% ------------------------------------------------------------------------------
\hypersetup{
    colorlinks=true,
    linkcolor=blue!80!black,
    citecolor=blue!80!black,
    urlcolor=blue!80!black
}

% ------------------------------------------------------------------------------
% Entornos Estilizados (Encabezado propio con título blanco y texto en la línea siguiente)
% ------------------------------------------------------------------------------

% 1. DEFINICIÓN (Azul)
\newtcbtheorem[number within=chapter]{definicion}{Definición}{
    colback=blue!3!white,
    colframe=blue!75!black,
    coltitle=white,
    fonttitle=\bfseries,
    enhanced,
    separator sign={: },
    breakable
}{def}

% 2. TEOREMA (Rojo)
\newtcbtheorem[number within=chapter]{teorema}{Teorema}{
    colback=red!3!white,
    colframe=red!75!black,
    coltitle=white,
    fonttitle=\bfseries,
    enhanced,
    separator sign={: },
    breakable
}{teo}

% 3. PROPOSICIÓN (Naranja)
\newtcbtheorem[number within=chapter]{proposicion}{Proposición}{
    colback=orange!3!white,
    colframe=orange!75!black,
    coltitle=white,
    fonttitle=\bfseries,
    enhanced,
    separator sign={: },
    breakable
}{prop}

% 4. LEMA (Turquesa)
\newtcbtheorem[number within=chapter]{lema}{Lema}{
    colback=teal!3!white,
    colframe=teal!75!black,
    coltitle=white,
    fonttitle=\bfseries,
    enhanced,
    separator sign={: },
    breakable
}{lem}

% 5. EJEMPLO (Gris)
\newtcbtheorem[number within=chapter]{ejemplo}{Ejemplo}{
    colback=gray!4!white,
    colframe=gray!60!black,
    coltitle=white,
    fonttitle=\bfseries,
    enhanced,
    separator sign={: },
    breakable
}{ejem}

% 6. OBSERVACIÓN / NOTA (Subtítulo en blanco sobre banda oscura)
\newtcolorbox{observacion}[1][]{
    colback=gray!5!white,
    colframe=gray!60!black,
    coltitle=white,
    fonttitle=\bfseries,
    title={Observación \ifstrempty{#1}{}{(#1)}},
    enhanced,
    breakable
}

% 7. DEMOSTRACIÓN
\newenvironment{demostracion}[1][Demostración]{%
    \par\vspace{0.5em}\noindent
    \tcbset{
        enhanced,
        colback=gray!3!white,
        colframe=gray!50!black,
        boxrule=0pt,
        leftrule=3pt,
        arc=0pt,
        outer arc=0pt,
        top=6pt, bottom=6pt, left=10pt, right=10pt,
        breakable
    }%
    \begin{tcolorbox}%
    \textbf{#1.}\,%
}{%
    \hfill$\square$%
    \end{tcolorbox}\par\vspace{0.5em}%
}

% 8. PROBLEMA RESUELTO (Púrpura)
\newtcbtheorem[number within=chapter]{problema}{Problema Resuelto}{
    colback=purple!3!white,
    colframe=purple!75!black,
    coltitle=white,
    fonttitle=\bfseries,
    enhanced,
    separator sign={: },
    breakable
}{prob}

% ------------------------------------------------------------------------------
% Notación Matemática
% ------------------------------------------------------------------------------
\newcommand{\R}{\mathbb{R}}
\newcommand{\N}{\mathbb{N}}
\newcommand{\Z}{\mathbb{Z}}
\newcommand{\Q}{\mathbb{Q}}

\newcommand{\cA}{\mathcal{A}}
\newcommand{\cB}{\mathcal{B}}
\newcommand{\cF}{\mathcal{F}}
\newcommand{\cG}{\mathcal{G}}

\newcommand{\Pprob}{\mathbb{P}}
\newcommand{\Eesperanza}{\mathbb{E}}
\newcommand{\Vvarianza}{\mathrm{Var}}
\newcommand{\medida}{\mu}